\chapter{Entrevista: Ing Paola Aliendre}
\label{sec:entrevista01}
\noindent
\textbf{Fecha:}03-04-2024

\noindent
\textbf{¿Cuál considera usted que es el aspecto más complejo entre los 13 pasos, puntos o documentos que se deben presentar al subir el PGSST a la página? ¿Por qué? ¿O cuál demanda más tiempo para completarse?}\\
De todos los 13 pasos, puntos o documentos que se presentan cuando se sube el PSCT a la página, el más complicado es el punto uno, que describe todas las condiciones de la empresa. Este punto es crucial ya que contiene el soporte técnico de cómo implementar la seguridad, y además se han añadido otros complementos como la gestión de la ergonomía y la parte psicosocial.
\\
\noindent
\textbf{¿Cuánto tiempo, en términos generales, toma completar todo el programa? Por ejemplo, en días o meses.}\\
Dependiendo del tamaño de la empresa, la implementación del programa completo puede llevar desde tres meses hasta incluso seis meses en el caso de pequeñas empresas sin ningún sistema previo.
\\
\noindent
\textbf{¿Cuáles características le gustaría que el sistema a desarrollar posea?}\\
Se desea un programa que facilite la aplicación de los métodos de gestión de seguridad, especialmente para las pequeñas y medianas empresas, con características que agilicen el proceso y reduzcan los tiempos de desarrollo y de cumplimiento legal.
\\
\noindent
\textbf{¿Podría proporcionarme acceso a una base de datos de proyectos de salud y seguridad previamente realizados para utilizarlos como referencia? ¿O qué alternativas sugiere para obtener recursos similares?}\\
No es posible obtener acceso a una base de datos de proyectos de salud y seguridad ya hechos, pero se sugiere buscar en las bases de datos de las universidades, como la UMSA, donde se pueden encontrar proyectos de grado relacionados con seguridad y salud laboral.
\\
\noindent
\textbf{¿Dónde podría obtener estadísticas referente al PGSST en Bolivia?}\\
La información sobre el sistema de gestión de seguridad se puede encontrar consultando registros en el Ministerio del Trabajo, donde se registran las empresas formales y sus sistemas de gestión.

%------------------------------------------------------------------------------
%------------------------------------------------------------------------------
\noindent
\textbf{Fecha:}17-06-2024

\noindent
\textbf{¿Cuáles son las actividades principales que se realizan para implementar un Programa de Gestión de Seguridad y Salud en el Trabajo (PGSST)?}\\
Las actividades principales incluyen la identificación de riesgos mediante la IPER, estudios de higiene y monitoreo, elaboración del plan de emergencia, planificación y capacitaciones, y organización del comité mixto. Cada una de estas actividades tiene su importancia y debe ser desarrollada de acuerdo con la normativa vigente.
\\
\noindent
\textbf{¿Qué documentos se deben desarrollar como parte de estas actividades para un PGSST?}\\
Los documentos que se deben desarrollar como parte de estas actividades para un PGSST incluyen la matriz de identificación de peligros, evaluación y control de riesgos (IPER), el plan de seguridad y salud en el trabajo, los procedimientos operativos estándar, los registros de capacitación, y los informes de auditoría.
\\
\noindent
\textbf{¿Cuál es el porcentaje de tiempo dedicado a cada tarea en la implementación de un PGSST?}\\
El tiempo varía según la empresa y la actividad. La IPER es una de las actividades que más tiempo demanda, aproximadamente un 40\% del tiempo total del proyecto. El plan de emergencia y el comité mixto pueden tomar alrededor de un 20\% y 10\% respectivamente. Los monitoreos también representan un 10\% del tiempo total. El resto de tiempo se distribuye entre las demás tareas.
\\
\noindent
\textbf{¿Qué aspectos deben considerarse al determinar los costos de implementar un PGSST?}\\
Los costos varían dependiendo de la magnitud de la empresa, la cantidad de trabajadores, el estado actual de la empresa y el tiempo que demandará la implementación. Por ejemplo, en el caso de Repsol, se cobra por campo petrolero, alrededor de 10 mil dólares cada uno. Para una empresa pequeña o mediana, los costos pueden ser significativamente menores, posiblemente alrededor de 300 dólares si se utiliza un software adecuado para reducir los costos.
\\
\noindent
\textbf{¿Cómo se cotizan las actividades específicas dentro del PGSST?}\\
Las actividades específicas, como la IPER, pueden ser cotizadas por ítem. Por ejemplo, para una empresa particular, el servicio de consultoría para la IPER puede costar varios miles de dólares, dependiendo de la complejidad y el tiempo requerido. No existe un arancel fijo, y todo se basa en la oferta y demanda, así como en la experiencia y reputación del profesional.
\\
\noindent
\textbf{¿Cómo se diferencian los costos y tiempos de implementación entre diferentes tipos de empresas, como pymes y grandes corporaciones?}\\
Las empresas más grandes, como YPFB, tienen estructuras más complejas y pueden requerir más tiempo y recursos para la implementación. Las pymes, en cambio, pueden tener una implementación más sencilla y rápida, pero aún así, todo depende de su estructura organizacional y el alcance del PGSST que deseen implementar.
\\
\noindent
\textbf{¿Podría dar una referencia de costos para diferentes tamaños de empresas?}\\
Para una microempresa, los costos podrían ser mucho menores, ya que la estructura y la cantidad de trabajo son más pequeñas. Personalmente, no cobro menos de 500 dólares por cualquier tipo de asesoría. Para proyectos más grandes, como el de Repsol, los costos pueden estar entre 10 mil y 14 mil dólares por campo petrolero.
\\
\noindent
\textbf{Analizando el caso específico del PGSST de la Universidad Católica, ¿cuáles fueron los pasos iniciales para su desarrollo?}\\
Los pasos iniciales incluyeron reuniones con los involucrados, la creación de una agenda de trabajo conjunto y la asignación de un supervisor o fiscal del proyecto. Este proceso tomó alrededor de tres meses y requirió la participación activa de los responsables de cada área.
\\
\noindent
\textbf{¿Cuánto tiempo tardó cada una de las partes del PGSST en la Universidad Católica?}\\
La IPER tomó aproximadamente dos meses y medio. El comité mixto fue lo más complicado, ya que no existía previamente y se demoró alrededor de cuatro a seis meses en ser conformado. Los monitoreos fueron rápidos, aproximadamente una semana, y la implementación del plan de acción se realizó en paralelo con estas actividades.
\\
\noindent
\textbf{¿Qué aceptación tuvo el PGSST en la Universidad Católica y qué cambios se notaron?}\\
La aceptación fue positiva, especialmente en términos de organización interna. Se nombró un comité mixto y se implementaron varias mejoras basadas en el programa. Sin embargo, el éxito a largo plazo depende de la implementación y seguimiento continuo por parte de la universidad.
\\
\noindent
\textbf{En su experiencia, ¿qué factores contribuyen a la falta de cumplimiento o implementación inadecuada de los PGSST?}\\
Los factores que contribuyen a la falta de cumplimiento o implementación inadecuada de los PGSST son la falta de compromiso por parte de la alta dirección, la insuficiencia de recursos financieros y humanos, y la carencia de una cultura de seguridad en la organización.
\\
\noindent
\textbf{¿Cómo se maneja la actualización y adaptación de los PGSST ante cambios en la legislación o normativas de seguridad y salud ocupacional?}\\
La actualización y adaptación de los PGSST ante cambios en la legislación o normativas de seguridad y salud ocupacional se maneja revisando periódicamente las políticas y procedimientos del programa, asegurando que cumplan con las nuevas regulaciones, y capacitando al personal sobre los cambios relevantes. Tu proyecto debería estar dirigido para aquellas empresas que se quieren implementar o actualizar su PGSST.
\\
\noindent
\textbf{¿Qué estrategias se han implementado para involucrar a los empleados en el proceso de desarrollo y mejora continua del PGSST?}\\
Las estrategias para involucrar a los empleados en el proceso de desarrollo y mejora continua del PGSST incluyen la creación de comités de seguridad, la realización de encuestas de opinión, y la implementación de programas de incentivos por cumplimiento de normas de seguridad.
\\
\noindent
\textbf{¿Cómo se puede fomentar una cultura de prevención y seguridad más proactiva dentro de las organizaciones?}\\
Para fomentar una cultura de prevención y seguridad más proactiva dentro de las organizaciones, es esencial implementar programas de capacitación continua, promover la participación activa de todos los niveles de la organización, y establecer sistemas de reconocimiento y recompensas para el cumplimiento de prácticas seguras.
\\
\noindent
\textbf{¿Cómo se utilizan los equipos para los estudio y monitoreos de higiene?}\\
En la práctica los equipos ya tienen softwares que directamente se conectan. Donde tú descargas los datos, pero lo difícil de eso es que igual tienes los datos y tienes que pasar a la planilla.

%------------------------------------------------------------------------------
%------------------------------------------------------------------------------
\noindent
\textbf{Fecha:}18-06-2024

\noindent
\textbf{¿Que actividades son las más satisfactorias profesionalmente del desarrollo de un PGSST?}\\
Las actividades más satisfactorias profesionalmente del desarrollo de un PGSST incluyen la implementación de las medidas de control. El entrevistado menciona que es gratificante ver cómo las medidas propuestas, basadas en la jerarquía de controles (eliminación, sustitución, controles de ingeniería, administrativos y EPP), se implementan y contribuyen a mejorar la seguridad y salud en el trabajo. Esta actividad requiere tanto de conocimientos teóricos como prácticos, y ver los resultados positivos de estas medidas proporciona una gran satisfacción profesional.
\\
\noindent
\textbf{¿Que actividades son las menos satisfactorias profesionalmente del desarrollo de un PGSST?}\\
Las actividades menos satisfactorias incluyen la investigación de accidentes, ya que se trata de analizar eventos después de que han ocurrido, lo cual no puede revertir los daños sufridos. También se mencionan los informes de monitoreo, que a menudo no conducen a acciones concretas si no se actúa sobre las recomendaciones.
%\\
%\noindent
%\textbf{¿Cuanto tiempo se invierte en introducir datos a plantillas o archivos de Excel?}\\
%Se estima que alrededor del 15-20\% del tiempo total del desarrollo de un PGSST se invierte en tareas administrativas, como la introducción de datos en plantillas o archivos de Excel, así como en la capacitación necesaria para completar estos registros adecuadamente.
\\
\noindent
\textbf{¿Qué actividades del desarrollo de una IPER son las mas y las menos satisfactorias profesionalmente?}\\
En el desarrollo de una IPER (Identificación de Peligros y Evaluación de Riesgos), las actividades más satisfactorias incluyen la identificación y cuantificación de riesgos y la implementación de medidas de control efectivas. Estas actividades permiten visualizar y mitigar los peligros de manera proactiva. Las menos satisfactorias son las relacionadas con peligros recurrentes y aparentemente "tontos" que, aunque pueden causar accidentes graves, a menudo son subestimados y requieren atención constante.
\\
\noindent
\textbf{¿Qué actividades del desarrollo de un plan de emergencia son las mas y las menos satisfactorias profesionalmente?}\\
En el desarrollo de un plan de emergencia, las actividades más satisfactorias incluyen la planificación de escenarios adecuados para la organización y la implementación de medidas prácticas y alcanzables que aseguren una respuesta eficaz en situaciones de emergencia. Las menos satisfactorias son aquellas en las que, a pesar de una buena planificación, la falta de aplicación práctica y de formación adecuada en la empresa reduce la efectividad del plan.
\\
\noindent
\textbf{¿Que actividades del desarrollo de un PGSST deberían automatizarse a su criterio?}\\
Basado en mi experiencia, hay varias actividades dentro del desarrollo de un Programa de Gestión de Seguridad y Salud en el Trabajo (PGSST) que deberían automatizarse para mejorar la eficiencia y reducir la carga administrativa sobre los profesionales. Una de las principales actividades que debería automatizarse es la gestión de los informes de monitoreo. Estos informes, aunque esenciales, a menudo se quedan en papel y no se implementan adecuadamente. Automatizar la recopilación y el análisis de los datos de monitoreo permitiría una respuesta más rápida y efectiva a las necesidades identificadas, como ajustes en la iluminación o reducción del ruido en el ambiente laboral.\\
Otra área que se beneficiaría de la automatización es la capacitación y la formación del personal en el llenado de registros y plantillas. Actualmente, se invierte mucho tiempo en enseñar a los trabajadores cómo completar estos documentos correctamente.\\
La implementación de políticas y procedimientos estándar también podría ser automatizada. Aunque cada empresa tiene sus particularidades, los fundamentos de las políticas de seguridad y salud ocupacional son similares. Un sistema que permita personalizar plantillas de políticas y procedimientos según los requisitos específicos de la empresa, pero basado en estándares como la ISO 45001, podría facilitar mucho este proceso.\\
Además, la automatización del diseño y la gestión de la matriz IPER (Identificación de Peligros y Evaluación de Riesgos) sería muy beneficiosa. La IPER es una herramienta fundamental que requiere tanto análisis cualitativo como cuantitativo. Un software que ayude a identificar peligros y evaluar riesgos basándose en datos históricos y en tiempo real mejoraría la precisión y reduciría el tiempo necesario para realizar estas evaluaciones.

%------------------------------------------------------------------------------
%------------------------------------------------------------------------------
\begin{center}
    \vspace{1.5cm}
    \hspace*{0pt}
    \textcolor{gray}{\rule{180pt}{0.1pt}}
    \vspace{-0.2cm}\\
    \MakeUppercase{\bf{ Paola Gioconda Aliendre Martinez }}\\[-0.2cm]
    \MakeUppercase{\bf{ Ing. Industrial, experta en Seguridad y Salud en el Trabajo }} 
\end{center}